\documentclass[12pt,table,xcolor={dvipsnames}]{beamer}
\usetheme{Pittsburgh}
\usecolortheme{seagull}
%\usepackage[utf8]{inputenc}
\usepackage{fontspec}
\usepackage[brazilian]{babel}
\usepackage{amsmath}
\usepackage{listings}
\usepackage{multirow}
\usepackage{amsfonts}
\usepackage{amssymb}
\usepackage{lstlinebgrd}
\usepackage{graphicx}
\author{Fila Encadeada, Pilha Encadeada e Lista Duplamente Encadeada}
\title{Estruturas de Dados}
%\setbeamercovered{transparent}
\setbeamertemplate{navigation symbols}{}
%\logo{\includegraphics[scale=0.015]{Brasao_UFSC.png}\includegraphics[scale=0.2]{brasao_PPGCC.jpg}}
\institute{Departamento de Informática e de Estatística \\ Prof. Jean Everson Martina \\ Prof. Aldo von Wangenheim}
\date{2016.2}
\subject{}
\usebackgroundtemplate{\includegraphics[width=\paperwidth,
height=\paperheight]{../reusable_images/fundo_UFSC.png}}
\begin{document}

{
\usebackgroundtemplate{\includegraphics[width=\paperwidth,
height=\paperheight]{../reusable_images/fundo_capa.png}}
\begin{frame}
\titlepage
\includegraphics[scale=0.3]{../reusable_images/brasao_INE.png}
\end{frame}
}

\begin{frame}[fragile]{Definições de Lista}
\begin{itemize}
\item 
\end{itemize}
\end{frame}


\begin{frame}[fragile]{Modelagem da Lista de Vetor}
\lstset{language=C++,
          keywordstyle=\color{blue}\ttfamily,
          stringstyle=\color{red}\ttfamily,
          commentstyle=\color{OliveGreen}\ttfamily,
          breaklines=true,
          basicstyle=\ttfamily\footnotesize
          }

\begin{itemize}
\item Aspecto Estrutural:
\begin{itemize}
\item Necessitamos de um vetor para armazenar as informações;
\item Necessitamos de um indicador da posição atual do ultimo elemento da lista;
\end{itemize}
\begin{lstlisting}
constantes MAXLISTA = 100;

classe Lista {
 T _dados[MAXLISTA]; // Vetor Estatico
 inteiro _ultimo;
};
\end{lstlisting}
\end{itemize}
\end{frame}



\begin{frame}[fragile]{Trabalho Lista de Vetor}
\lstset{language=C++,
          keywordstyle=\color{blue}\ttfamily,
          stringstyle=\color{red}\ttfamily,
          commentstyle=\color{OliveGreen}\ttfamily,
          breaklines=true,
          basicstyle=\ttfamily\footnotesize
          }
\begin{itemize}
\item Implemente uma classe Lista todas as operações vistas;
\item Implemente a lista usando Templates;
\item Implemente a lista com um numero de elementos variável definido na instanciação;
\item Use as melhores práticas de orientação a objetos;
\item Documente todas as classes, métodos e atributos;
\item Aplique os testes unitários disponíveis no moodle da disciplina para validar sua estrutura de dados;
\item Entregue até a data definida no moodle.
\end{itemize}
\end{frame}

{
\usebackgroundtemplate{\includegraphics[width=\paperwidth,
height=\paperheight]{../reusable_images/fundo_capa.png}}
\begin{frame}

{\LARGE Perguntas????}

\end{frame}
}


{
\usebackgroundtemplate{\includegraphics[width=\paperwidth,
height=\paperheight]{../reusable_images/fundo_capa.png}}
\begin{frame}
\includegraphics[scale=0.8]{../reusable_images/cc_logo_arge.png}\hspace{0.5cm}
\includegraphics[scale=0.95]{../reusable_images/by.png}

\vspace{1cm}
Este trabalho está licenciado sob uma Licença Creative Commons Atribuição 4.0 Internacional. Para ver uma cópia desta licença, visite http://creativecommons.org/licenses/by/4.0/.

\end{frame}
}
\end{document}
